\documentclass[11pt]{article}
\usepackage[margin=26mm]{geometry}
\usepackage{amsmath,amssymb,amsthm}
\newtheorem{proposition}{Proposition}
\title{A Four-Vertex Counterexample to a Spanning-Path Statement\\\large Known Control for Conjecture Harness}
\author{Software demonstration; no research author attributed}
\date{}
\begin{document}
\maketitle
\begin{abstract}
We document the elementary obstruction to a Hamiltonian path in the four-vertex star. This known control demonstrates a reproducibility workflow; no novelty is claimed.
\end{abstract}
\section{Statement and definitions}
Consider the false statement that every finite connected simple graph has a Hamiltonian path. A Hamiltonian path is a sequence of distinct vertices containing all vertices of the graph, with an edge between each consecutive pair.
\begin{proposition}
The star $K_{1,3}$ is a connected simple graph with no Hamiltonian path.
\end{proposition}
\begin{proof}
Let the vertex set be $\{0,1,2,3\}$ and the edge set be $\{\{0,1\},\{0,2\},\{0,3\}\}$. The graph is simple and connected because every other vertex is adjacent to $0$. Each of $1,2,3$ has degree one. An internal vertex of a simple path must have two distinct neighbors on that path. Consequently all three leaves must be endpoints of any Hamiltonian path. A path on four vertices has exactly two endpoints, a contradiction.
\end{proof}
\section{Reproducibility and scope}
This demonstration uses the development version of Conjecture Harness~\cite{conjectureharness2026}. The accompanying Python script reads the exact edge list, checks connectivity and the leaf-degree obstruction, and enumerates all $4!=24$ vertex permutations. None is a Hamiltonian path. Enumeration checks the witness only, not a general theorem. Run \texttt{python repro/verify.py} from the package root. The output is included in the reproduction directory.
\section*{Acknowledgments and AI disclosure}
ChatGPT generated this known-control demonstration and accompanying code. No human research author, independent peer review or priority claim is asserted.
\begin{thebibliography}{99}
\raggedright
\bibitem{conjectureharness2026} KAVentures. Conjecture Harness: A Chat-First Protocol for Reproducible AI-Assisted Mathematics. 2026. Version 0.4.0-dev (unreleased local snapshot). https://github.com/KAVentures/conjecture-harness.
\end{thebibliography}

\end{document}
